\documentclass[11pt, letterpaper]{article}
\usepackage[utf8]{inputenc}
\usepackage{geometry}
\geometry{margin=1in}
\usepackage{amsmath}
\usepackage{amssymb}
\usepackage{graphicx}
\usepackage{hyperref}

\title{\vspace{-2cm}\textbf{QKD Optimization Performance Analysis:\\ The 3,574$\times$ Hybrid Speedup}}
\author{Systems Intelligence \& Algorithmic Efficiency}
\date{}

\begin{document}

\maketitle
\vspace{-1.5cm}

\section*{1. Simulation Domain \& Baseline Assumptions}

Let the spatial simulation sequence for the QKD optimization be defined over $N_p$ discrete points (from $0$ to $200$ km with $\Delta L = 0.5$ km):
\begin{equation}
N_p = 400 \text{ points}
\end{equation}

For the legacy stochastic solver (Dual Annealing), the number of objective function evaluations required to reach the global minimum per spatial point is:
\begin{equation}
N_e \approx 15,000 \text{ evaluations/point}
\end{equation}


\section*{2. Legacy Baseline: Standard Python + Stochastic Search}

Let $t_{\text{Py}}$ be the raw evaluation time of the Lim et al. 2014 finite-key decoy-state bounds in an interpreted Python environment:
\begin{equation}
t_{\text{Py}} \approx 47 \text{ ms} = 0.047 \text{ s}
\end{equation}

The compute time required for a single spatial point, $T_{\text{point, base}}$, and the total sequence time, $T_{\text{total, base}}$, are:
\begin{align}
T_{\text{point, base}} &= N_e \times t_{\text{Py}} \approx 15,000 \times 0.047 \text{ s} = 705 \text{ s} \\
T_{\text{total, base}} &= T_{\text{point, base}} \times N_p = 705 \text{ s} \times 400 = 282,000 \text{ s} \approx 78.3 \text{ Hours}
\end{align}


\section*{3. Breakthrough 1: Systems Engineering (JAX + XLA Compilation)}

By compiling the physics engine into XLA machine code, the evaluation time drops to $t_{\text{JAX}}$:
\begin{equation}
t_{\text{JAX}} \approx 102 \ \mu\text{s} = 1.02 \times 10^{-4} \text{ s}
\end{equation}

The new computational time for a single spatial point becomes:
\begin{equation}
T_{\text{point, JAX}} = N_e \times t_{\text{JAX}} \approx 15,000 \times 1.02 \times 10^{-4} \text{ s} = 1.53 \text{ s}
\end{equation}

The total sequence time (including compilation overhead and Python loops) is empirically measured as $T_{\text{total, JAX}} = 709.0 \text{ s}$.

The isolated Systems Speedup factor ($S_{\text{sys}}$) at the raw objective level is:
\begin{equation}
S_{\text{sys}} = \frac{T_{\text{point, base}}}{T_{\text{point, JAX}}} = \frac{705 \text{ s}}{1.53 \text{ s}} \approx 460\times
\end{equation}
\textit{(Conservatively reported as $>100\times$ on slides to account for global hardware variance).}


\section*{4. Breakthrough 2: Algorithmic Efficiency (Exact FoR Hessian)}

By replacing the $N_e \approx 15,000$ stochastic guesses with an exact Forward-over-Reverse 2nd-order Hessian tracker (\texttt{trust-constr}), the number of required calculations per point collapses. The empirically measured time to compute the entire 400-point sequence drops to $T_{\text{total, Hess}} = 78.9 \text{ s}$.

The isolated Algorithmic Speedup factor ($S_{\text{alg}}$) is:
\begin{equation}
S_{\text{alg}} = \frac{T_{\text{total, JAX}}}{T_{\text{total, Hess}}} = \frac{709.0 \text{ s}}{78.9 \text{ s}} \approx 8.99\times
\end{equation}


\section*{5. The Compounding Total Speedup ($S_{\text{total}}$)}

The overall engineering achievement combines the execution speed of the JAX compiler ($S_{\text{sys}}$) with the search efficiency of the exact Hessian ($S_{\text{alg}}$). The total computational acceleration against the standard legacy baseline is:
\begin{equation}
S_{\text{total}} = \frac{T_{\text{total, base}}}{T_{\text{total, Hess}}} = \frac{282,000 \text{ s}}{78.9 \text{ s}} \approx \textbf{3,574}\times
\end{equation}

\vspace{0.5cm}
\hrule
\vspace{0.3cm}
\noindent \small \textit{Conclusion:} A computationally intractable analysis (78.3 hours) was reduced to an interactive timescale (78.9 seconds) strictly through rigorous application of Systems Engineering and Deep Mathematics.

\end{document}
